\documentclass[11pt]{article}

% Recommended, but optional, packages for figures and better typesetting:
\usepackage{microtype}
\usepackage{graphicx}
\usepackage{subfigure}
\usepackage{booktabs} % for professional tables

% hyperref makes hyperlinks in the resulting PDF.
\usepackage[colorlinks]{hyperref}

% Attempt to make hyperref and algorithmic work together better:
\newcommand{\theHalgorithm}{\arabic{algorithm}}
\usepackage{algorithm}
\usepackage{algorithmic}

% --- Handle long URLs in bibliography ---
\usepackage{url}
\expandafter\def\expandafter\UrlBreaks\expandafter{\UrlBreaks\do\/\do\-\do\.\do\0\do\1\do\2\do\3\do\4\do\5\do\6\do\7\do\8\do\9\do\a\do\b\do\c\do\d\do\e\do\f\do\g\do\h\do\i\do\j\do\k\do\l\do\m\do\n\do\o\do\p\do\q\do\r\do\s\do\t\do\u\do\v\do\w\do\x\do\y\do\z\do\A\do\B\do\C\do\D\do\E\do\F\do\G\do\H\do\I\do\J\do\K\do\L\do\M\do\N\do\O\do\P\do\Q\do\R\do\S\do\T\do\U\do\V\do\W\do\X\do\Y\do\Z}

% Use the following line for the initial blind version submitted for review:
\usepackage{acl}

\usepackage{amsmath}
\usepackage{amssymb}
\usepackage{mathtools}
\usepackage{amsthm}

% --- Custom Packages ---
\usepackage{multirow}
\usepackage{enumitem}
\usepackage{siunitx}
\usepackage{listings}
\usepackage{xcolor}

% --- Chinese Support (pdflatex compatible) ---
\usepackage{CJKutf8}

\sisetup{detect-all}
\usepackage[capitalize,noabbrev]{cleveref}

% --- Code Block Styles ---
\lstdefinelanguage{json}{
  basicstyle=\ttfamily\small,
  numbers=left,
  numberstyle=\footnotesize\color{gray},
  stepnumber=1,
  numbersep=6pt,
  showstringspaces=false,
  breaklines=true,
  frame=single,
  rulecolor=\color{black!20},
  string=[s]{"}{"},
  comment=[l]{//},
  morecomment=[s]{/*}{*/},
  keywordstyle=\color{blue!60!black},
  stringstyle=\color{green!50!black},
}

\lstdefinestyle{prompt}{
  basicstyle=\ttfamily\small,
  numbers=left,
  numberstyle=\footnotesize\color{gray},
  stepnumber=1,
  numbersep=6pt,
  showstringspaces=false,
  breaklines=true,
  frame=single,
  rulecolor=\color{black!20},
}

\title{MedTime: A Differentiable Consistency Framework for \\ Temporal Reasoning over Clinical Narratives}

\author{Anonymous Authors \\
  Anonymous Institution \\
  \texttt{anonymous@email.invalid} \\}

\begin{document}
\begin{CJK*}{UTF8}{gbsn}
\maketitle

\begin{abstract}
Extracting consistent patient timelines from clinical narratives is essential for disease trajectory modeling and prognosis prediction. However, current methods often treat temporal cues in isolation, leading to global inconsistencies across fragmented or multi-document records. We propose \textbf{MedTime}, a framework that reformulates temporal reasoning as a differentiable consistency problem. MedTime unifies schema-constrained generation, retrieval-augmented contextual verification, and a differentiable projection layer to enforce transitive coherence during model optimization. Built on \texttt{LLaMA-Med-13B}, our approach enables end-to-end learning of consistent timelines without the need for external post-hoc solvers. Evaluations on the \textbf{MedAlign-CN} benchmark show that MedTime improves temporal F1 by 5.6 points and reduces timeline edit distance by 32\% over strong baselines, demonstrating that differentiable structural regularization is critical for resolving implicit temporal dependencies in clinical text.
\end{abstract}

\section{Introduction}

Reliable temporal reasoning is a foundational requirement for clinical applications such as disease trajectory modeling and time-aware patient characterization. Despite the maturity of general information extraction (IE), clinical temporal reasoning remains a significant challenge due to the implicit nature of temporal anchors, fragmented documentation, and the scarcity of high-quality annotated data. Early work focused on interval algebra \cite{allen1983maintaining} and relation classification using schemes like MATRES \cite{ning2019matres}, while more recent neural models \cite{lin2020joint,wang2022mlee} have improved span-level extraction. However, most existing systems treat temporal links as independent pairwise decisions, lacking a mechanism to maintain global coherence across an entire patient episode.

Current temporal IE approaches typically process pairwise temporal links (\textsc{TLINKs}) \cite{han2023clinicaltimeline,lin2023temptime}. In clinical narratives, where temporal references are often cross-document or underspecified \cite{lin2023temptime,su2025transformerTIE}, local pairwise models frequently predict relations that are locally plausible but globally contradictory. While post-hoc closure algorithms \cite{ning2019matres,han2023clinicaltimeline} can enforce transitive consistency, they do not resolve underlying ambiguities or provide a gradients to improve the model's internal representations during training. Even state-of-the-art LLMs with retrieval augmentation \cite{li2024gptmed,wang2024medtemporal,fan2025consistent,eirew2025beyondpairwise,zhao2025temporaleval} struggle to respect complex structural priors like transitivity and acyclicity in multi-document settings.

To address these limitations, we introduce \textbf{MedTime}, a \emph{diferentiable Generate--Verify--Project (GVP) framework} that integrates event extraction and temporal alignment under global consistency constraints. Unlike pipeline-based repair, MedTime incorporates temporal logic directly into the objective function through three synergistic components: a schema-constrained generator for structured event-graph production; a contextual verifier to resolve implicit or cross-document evidence; and a differentiable projection layer that enforces transitivity and drifting constraints. By optimizing these modules end-to-end, MedTime transforms discrete temporal reasoning into a differentiable process, allowing the model to internalize structural validity during generation.

We evaluate MedTime on \textbf{MedAlign-CN}, a new Chinese clinical benchmark designed for multi-document temporal reasoning. Our findings show that MedTime achieves significant gains, improving temporal F1 by 5.6 points and reducing timeline edit distance by 32\% compared to previous state-of-the-art models. Our contributions are three-fold:
\begin{enumerate}[nosep]
  \item We introduce the \emph{Generate--Verify--Project} paradigm, providing a differentiable alternative to independent pair classification and post-hoc repair.
  \item We present a \emph{differentiable temporal projection} layer that embeds logical transitivity and temporal drift constraints directly into the training objective.
  \item We provide empirical evidence on MedAlign-CN demonstrating that soft structural regularization leads to more coherent and reliable clinical timelines than standard RAG-based LLM architectures.
\end{enumerate}

\section{Related Work}
\label{sec:related-work}

\subsection{Span- and MRC-based Event Extraction}
Early biomedical IE research typically formulated event extraction as span classification or machine reading comprehension (MRC), identifying event triggers and arguments independently~\cite{lin2020joint,wang2022mlee}. Such systems achieve strong local precision but treat each mention in isolation, without reasoning over temporal or causal order. These span-based formulations constitute our \emph{local extraction baselines} in~\S\ref{sec:exp}, and motivate the need for structured representations that capture both event semantics and temporal coherence.

\subsection{Temporal IE and Global Consistency}
Temporal IE has evolved from interval algebra~\cite{allen1983maintaining} to neural architectures predicting temporal links (TLINKs) between event pairs. Recent models—TempTime~\cite{lin2023temptime}, TIE-Transformer~\cite{su2025transformerTIE}, Clinical-Timeline~\cite{han2023clinicaltimeline}, and span-based graph transformers~\cite{chaturvedi2025temporal}—extend reasoning from sentence to document level, while cross-lingual and zero-shot studies~\cite{kougia2024zeroshotclinicalTRE} highlight the difficulty of maintaining temporal consistency across domains. However, these systems still rely on discrete closure algorithms applied after inference. Such post-hoc reasoning cannot guarantee global transitivity or bounded drift, and consistency is not optimized jointly with model parameters. Our work differs by introducing a \emph{differentiable temporal projection} that integrates global consistency directly into training, bridging symbolic temporal logic and neural optimization.

\subsection{LLMs, RAG, and Structured Temporal Reasoning}
Large language models (LLMs) have shown strong potential in clinical information extraction and temporal inference~\cite{li2024gptmed,wang2024medtemporal,lopez2025clear,fan2025consistent,eirew2025beyondpairwise}. Instruction-tuned or retrieval-augmented (RAG) LLMs can perform end-to-end temporal reasoning, yet they often generate non-transitive or contradictory timelines \cite{lu2023structureddecoding,welleck2024consistency,zhao2025temporaleval,kumar2025temporalbench,gong2025eventrelbench}. To mitigate this, prior work explored self-consistency and structured decoding strategies, but these methods still operate at the output level, without an explicit mechanism to optimize consistency constraints during training. \textbf{MedTime} extends this line by coupling schema-constrained generation, a contextual verifier that leverages retrieved evidence, and a \emph{differentiable temporal projection layer} that enforces transitivity and bounded drift in a unified optimization process.

\subsection{Neuro-Symbolic Reasoning in Medicine}
Recent trends in medical AI emphasize neuro-symbolic approaches that combine neural flexibility with logical guarantees. While standard LLMs often hallucinate or contradict themselves, frameworks integrating symbolic constraints (e.e., logic rules, knowledge graphs) have shown promise in ensuring reliability. Our work aligns with this direction by embedding temporal logic (transitivity) as a differentiable objective, effectively teaching the neural model to respect the physical laws of time.

\subsection{Summary}
Prior research spans three directions: (1) span/MRC-based local extraction, (2) pairwise temporal IE with post-hoc closure, and (3) LLM-based or RAG-enhanced structured reasoning. While each tackles parts of the temporal reasoning challenge, none unifies event generation, verification, and consistency enforcement within a single differentiable end-to-end framework. \textbf{MedTime} bridges these efforts through a \emph{structure-aware generate–Verify–project paradigm} that performs event extraction, temporal alignment, and global consistency enforcement jointly.

\section{Task and Problem Formulation}
\label{sec:task}

\paragraph{Motivation.}
Clinical narratives describe a patient's clinical course as a sequence of temporally anchored events (e.g., diagnosis, medication, surgery). However, these events are often dispersed across heterogeneous notes and expressed through implicit or underspecified temporal cues. Our goal is to extract and align such events into a \textbf{globally consistent event--time graph} that preserves both causal and temporal ordering across documents.

\paragraph{Event Definition.}
An event is defined as $e=\langle t,\mathcal{A},\tau\rangle$, where $t$ denotes the trigger, $\mathcal{A}=\{a_1,\dots,a_m\}$ are arguments (e.g., agent, target, value), and $\tau$ is its temporal anchor (e.g., an absolute or relative time expression). Given a collection of clinical documents for a patient, the model generates a structured event set $\mathcal{E}=\{e_1,\dots,e_n\}$ covering all salient events.

\paragraph{Temporal Relations and Global Consistency.}
Pairs of events $(e_i,e_j)$ are linked by temporal relations $\mathcal{R}_{ij}\!\in\!\{\textsc{Before},\textsc{Overlap},\textsc{After}\}$, forming a patient-level temporal graph $\mathcal{G}=(\mathcal{E},\mathcal{R})$. We use $e_i\!\prec\!e_j$ to denote that event $e_i$ occurs before $e_j$. A valid temporal graph must satisfy \textbf{global temporal consistency}:
\begin{equation}
\label{eq:consistency}
(e_i\!\prec\!e_j)\!\land\!(e_j\!\prec\!e_k) \Rightarrow(e_i\!\prec\!e_k), \qquad |\tau_i-\tau_j|\le\delta,
\end{equation}
where $\tau_i$ is the normalized anchor of $e_i$ and $\delta$ bounds the maximum allowable temporal drift. These constraints ensure transitivity, prevent cycles, and regularize long-range ordering across documents.

\paragraph{Learning Objective.}
Given a patient-level text collection $\mathcal{X}$, the model learns to generate a structured event--time graph $\hat{\mathcal{G}}=(\hat{\mathcal{E}},\hat{\mathcal{R}})$ that jointly satisfies extraction and ordering constraints. The overall loss integrates three components:
\begin{equation}
\label{eq:train_total}
\mathcal{L} =\mathcal{L}_{\text{event}} +\mathcal{L}_{\text{time}} +\lambda\,\mathcal{L}_{\text{consistency}},
\end{equation}
where $\mathcal{L}_{\text{event}}$ supervises trigger and argument generation from the LLM decoder, $\mathcal{L}_{\text{time}}$ aligns normalized temporal anchors, and $\mathcal{L}_{\text{consistency}}$ penalizes violations of Eq.~(\ref{eq:consistency}) through a differentiable projection layer. This formulation departs from post-hoc rule closure by optimizing global temporal validity \emph{during training}—the key design that enables MedTime’s generate–Verify–project framework to remain fully differentiable.

\paragraph{Output Representation.}
The final output is serialized as a structured event--time JSON graph (Fig.~\ref{fig:json}), facilitating downstream reasoning and visualization of patient trajectories. Each node corresponds to an extracted event, and each edge represents a consistent temporal relation inferred under the global constraint.

\begin{figure}[t]
\centering
\footnotesize
\begin{verbatim}
{
  "patient_id": "P001",
  "events": [
    {"trigger": "surgery",
     "arguments": {"anatomy": "liver"},
     "time": "2022-05-06",
     "relation": {"chemotherapy": "Before"}}
  ]
}
\end{verbatim}
\caption{Simplified example of structured event--time JSON output.}
\label{fig:json}
\end{figure}

\section{Method: Structure-Aware GVP Framework}
\label{sec:method}

We introduce \textsc{MedTime}, utilizing a \emph{generate–Verify–project} (GVP) architecture to unify extraction with logical consistency. Leveraging parameter-efficient fine-tuning via LoRA \cite{hu2021lora} on medical backbones \cite{wu2023pmcllama}, our framework builds on neuro-symbolic principles \cite{li2023scallop} and decomposed reasoning \cite{khot2023decomposed}. Unlike pipeline approaches that fix errors after the fact, our objective function (Eq.~\ref{eq:train_total}) trains three differentiable modules in unison: (1) \textbf{Generate} — producing the raw event graph; (2) \textbf{Verify} — retrieving evidence for ambiguous cues; and (3) \textbf{Project} — applying soft constraints for transitivity. This design ensures gradients propagate from consistency violations back to the generator.

\subsection{Generate: Schema-Constrained Decoder}
\textit{Goal:} Generate structured event–time graphs with valid schema and continuous temporal anchors.

Clinical narratives contain nested and cross-sentence dependencies, making sequence tagging inadequate. We therefore model event extraction as constrained text generation with a medical LLM backbone (\texttt{LLaMA-Med-13B} with \textsc{LoRA} adapters):
\begin{equation}
P(Y|X) = \prod_{t=1}^{T} P(y_t \mid y_{<t}, X),
\label{eq:decoder}
\end{equation}
where $X$ is the input text and $Y$ is a serialized JSON instance representing the patient-level event–time graph.

\paragraph{Schema-Constrained Decoding.}
Decoding follows a task-specific JSON schema (Appendix~\ref{app:prompts}) requiring valid \texttt{trigger}, \texttt{arguments}, and \texttt{time} fields. Constrained beam search with dynamic pruning discards malformed generations, ensuring structural validity and type completeness. This yields the \emph{event loss}:
\begin{equation}
\mathcal{L}_{\text{event}} = -\sum_{t=1}^{T} \log P(y_t^\ast \mid y_{<t}, X),
\label{eq:event_loss}
\end{equation}
which supervises token-level decoding including triggers and arguments. For each predicted event $e_i$, we obtain its contextual embedding $\mathbf{h}_i$ by mean-pooling token representations across its trigger span. These continuous vectors serve as the basis for differentiable temporal reasoning.

\paragraph{Soft Temporal Anchors.}
Each event’s time expression—absolute (e.g., ``2022-05-06'') or relative (e.g., ``postoperative day 3'')—is mapped to a continuous latent anchor $\hat{\tau}_i \in \mathbb{R}$ using a linear projection from the LLM hidden state corresponding to the \texttt{time} field:
\begin{equation}
\hat{\tau}_i = \mathbf{w}_\tau^\top \mathbf{h}_i + b_\tau.
\label{eq:tau_estimation}
\end{equation}
Unlike discrete normalization, $\hat{\tau}_i$ remains differentiable, so gradients from later consistency losses can flow back into the generator.

\subsection{Verify: Learned Contextual Verifier}
\textit{Goal:} Estimate the reliability of event pairs through logical calibrated scoring.

While embeddings $\mathbf{h}_i$ implicitly encode context from the LLM's RAG-augmented input, explicit confidence calibration is needed for the global projection. We introduce a lightweight \textbf{Learned Verifier} (an MLP head) that predicts the pairwise consistency weight $w_{ij}$ from the concatenated event representations:
\begin{equation}
w_{ij} = \sigma(\mathbf{W}_v [\mathbf{h}_i; \mathbf{h}_j] + b_v),
\label{eq:pair-simple}
\end{equation}
where $\sigma$ is the sigmoid function. This weight $w_{ij} \in [0,1]$ quantifies the model's distinct confidence in the relation between $e_i$ and $e_j$. It acts as a gate for the topological loss: we heavily penalize transitivity violations only for pairs where the model is confident in the link ($w_{ij} \to 1$), avoiding noise propagation from ambiguous mentions.

\paragraph{Calibration Objective.}
The verifier is jointly trained to predict valid relations (Eq.~\ref{eq:pair-simple}). The confidence $w_{ij}$ modulates the alignment loss $\mathcal{L}_{\text{time}}$:
\begin{equation}
\label{eq:time_loss}
\mathcal{L}_{\text{time}} = -\sum_{(i,j)} w_{ij}\, \log p(r_{ij} \mid \mathbf{h}_i, \mathbf{h}_j).
\end{equation}
By optimizing this objective, the verifier learns to down-weight uncertain pairs (where $\mathcal{L}_{\text{time}}$ would be high), effectively ignoring ambiguous timeline segments during global projection.

\subsection{Project: Differentiable Temporal Projection}
\textit{Goal:} Enforce global consistency through a Generate--Verify--Project (GVP) objective.

Clinical narratives often yield non-transitive timelines when predicted independently. To ensure global coherence, we project latent event representations into a differentiable constraint space. The total GVP loss is defined as:
\begin{equation}
\mathcal{L}_{\text{GVP}} = \alpha \mathcal{L}_{\text{align}} + \lambda \mathcal{L}_{\text{topo}} + \beta \mathcal{L}_{\text{ground}} + \gamma \mathcal{L}_{\text{drift}}.
\end{equation}

\paragraph{Alignment Loss ($\mathcal{L}_{\text{align}}$).}
We enforce semantic alignment between the event trigger embedding $\mathbf{h}_{trigger, i}$ and its projected temporal anchor $\mathbf{h}_{time, i}$ in the embedding space:
\begin{equation}
\mathcal{L}_{\text{align}} = \frac{1}{N} \sum_{i=1}^N \| \mathbf{h}_{trigger, i} - \mathbf{h}_{time, i} \|_2^2.
\end{equation}

\paragraph{Topological Loss ($\mathcal{L}_{\text{topo}}$).}
We utilize a differentiable logic head to penalize transitive violations. For event triplets $(e_i, e_j, e_k)$, the loss is:
\begin{equation}
\label{eq:topo_loss}
\begin{split}
\mathcal{L}_{\text{topo}} &= \sum_{i,j,k} w_{ijk}\,\text{softplus}(\hat{\tau}_j - \hat{\tau}_i \\ 
&\quad + \hat{\tau}_k - \hat{\tau}_j - (\hat{\tau}_k - \hat{\tau}_i) - m),
\end{split}
\end{equation}
where the violation term simplifies to $\text{softplus}(D_{ij} + D_{jk} - D_{ik} - m)$, enforcing the transitive property $D_{ik} \approx D_{ij} + D_{jk}$. $m=0.1$ serves as a margin.

\paragraph{Grounding Loss ($\mathcal{L}_{\text{ground}}$).}
To prevent semantic drift, we constrain the trigger embedding to remain close to the mean embedding of its retrieved contextual evidence $\mathcal{C}_i$:
\begin{equation}
\mathcal{L}_{\text{ground}} = \frac{1}{N} \sum_{i=1}^N \| \mathbf{h}_{trigger, i} - \text{mean}(\text{Enc}(\mathcal{C}_i)) \|_2^2.
\end{equation}

\paragraph{Bounded Drift Penalty ($\mathcal{L}_{\text{drift}}$).}
Finally, we penalize temporal intervals that exceed realistic clinical durations (e.g., massive gaps between post-op checks):
\begin{equation}
\mathcal{L}_{\text{drift}} = \sum_{i,j} \max(0, |D_{ij}| - \delta_{\text{max}}),
\end{equation}
where $\delta_{\text{max}}$ is the maximum allowable drift. Algorithm~\ref{alg:soft-proj-train} outlines the full differentiable projection process.

\subsection{Joint Optimization and Inference}
\textit{Goal:} Optimize generation, validation, and projection jointly for end-to-end consistency.

The total loss is $\mathcal{L} = \mathcal{L}_{\text{event}} + \mathcal{L}_{\text{time}} + \lambda\,\mathcal{L}_{\text{consistency}}$ (Eq.~\ref{eq:train_total}), where $\lambda$ balances local and global constraints. All three modules are trained jointly in a single differentiable computation graph. During inference, schema-constrained decoding ensures output validity, and a lightweight \emph{hard acyclic closure} (Algorithm~\ref{alg:hard-closure-test}) enforces residual transitivity, producing globally coherent timelines. This hybrid training–inference design bridges symbolic temporal logic and neural optimization, achieving interpretable, consistent, and end-to-end trainable clinical event reasoning.

\begin{algorithm}[t]
\small
\caption{Soft Transitive Projection for Training}
\label{alg:soft-proj-train}
\begin{algorithmic}[1]
\REQUIRE Soft anchors $\{\hat{\tau}_i\}$, weights $w_{ij}$, bounds $\delta_{ij}$
\STATE $D_{ij} \leftarrow \mathrm{clamp}(\hat{\tau}_j - \hat{\tau}_i,\,-C,\,C)$
\FOR{$p = 1$ to $P$}
  \STATE \textbf{Triangle:} $v_{ijk} = \mathrm{softplus}(D_{ij} + D_{jk} - D_{ik} - m)$
  \STATE \textbf{Accumulate:} $\mathcal{L}_{\text{tri}} += \sum \text{softmax}(w) v_{ijk}$
  \STATE \textbf{Drift:} $\mathcal{L}_{\text{bd}} += \sum w \text{softplus}(|D|-\delta)$
\ENDFOR
\STATE \textbf{return} $\mathcal{L}_{\text{consistency}}$
\end{algorithmic}
\end{algorithm}

\begin{algorithm}[t]
\small
\caption{Hard Acyclic Closure for Inference}
\label{alg:hard-closure-test}
\begin{algorithmic}[1]
\REQUIRE Predicted $D_{ij}$
\STATE $D \leftarrow \mathrm{clamp}(D,\,-C,\,C)$
\FOR{$p=1$ to $P$}
  \IF{$D_{ij}+D_{jk} < D_{ik}$}
     \STATE $D_{ik} \leftarrow D_{ij}+D_{jk}$
  \ENDIF
\ENDFOR
\STATE Threshold to labels
\end{algorithmic}
\end{algorithm}

\section{MedAlign-CN: Benchmark}
\label{sec:dataset}

We construct \textbf{MedAlign-CN}, the first Chinese benchmark explicitly designed to evaluate \emph{global temporal consistency} across heterogeneous clinical documents. Beyond serving as training data for \textsc{MedTime}, it provides a diagnostic suite for analyzing consistency violations in large language models.

\subsection{Sources and De-identification}
MedAlign-CN is curated from institutional electronic health records (\textit{admission notes, progress notes, imaging reports, discharge summaries}) collected between 2018–2023 under formal data-sharing agreements. De-identification combines deterministic regular expressions and transformer-based redaction. All records undergo manual PHI review before release. The corpus is distributed under a non-commercial research license.

\subsection{Annotation Schema and Protocol}
Each \textbf{patient episode} includes multiple documents and serves as one temporal reasoning unit. Annotators label \textbf{event triggers}, \textbf{arguments}, and \textbf{temporal anchors} (\textit{absolute} or \textit{relative}, e.g., “术后第3天” / postoperative day~3). For each event pair $(e_i, e_j)$, annotators assign a relation $r_{ij} \in \{\textsc{Before}, \textsc{Overlap}, \textsc{After}\}$. Cross-document links align diagnosis, treatment, and outcome events within the same patient. The annotation tool automatically computes temporal closure under Eq.~(\ref{eq:consistency}) and flags inconsistent cycles for adjudication.

\begin{figure}[t]
\centering
\fbox{
\begin{minipage}[c][40mm][c]{0.93\linewidth}
\small
\textbf{Example (CN → EN):}\\
\texttt{术后第 3 天复查CT，未见出血。}\\
\textit{CT re-examination on postoperative day~3 showed no hemorrhage.}\\[0.3em]
\textbf{Annotation:}\\
Trigger = “复查CT”; Anchor = “术后第3天”; Relation = \textsc{After(surgery)}.
\end{minipage}}
\caption{Example from MedAlign-CN illustrating event trigger, temporal anchor, and relation annotation.}
\label{fig:cn_example}
\vspace{-0.5em}
\end{figure}

\subsection{Quality Control and Agreement}
All documents are double-annotated by trained medical annotators and adjudicated by a senior clinician. Inter-annotator agreement (IAA) is measured on event spans, temporal links, and timeline consistency, following the evaluation protocol of \cite{lin2023temptime}. Final IAA reaches 0.91 (span~F1), 0.85 (TLINK~F1), and 0.82 (timeline consistency), surpassing the 0.8 reliability threshold commonly accepted in temporal IE.

\subsection{Corpus Scale and Statistics}
MedAlign-CN contains \textbf{150 patient episodes} ($\approx$2{,}000 documents) across surgery, oncology, and internal medicine. Each episode includes 6.8 documents and 95 annotated events on average, totalling \textbf{14.3k events} and \textbf{11.7k temporal links}.

\setlength{\tabcolsep}{4pt}
\begin{table}[t]
\centering
\resizebox{\columnwidth}{!}{
\begin{tabular}{lcccc}
\toprule
\textbf{Subset} & \textbf{Patients} & \textbf{Docs} & \textbf{Events} & \textbf{TLINK F1} \\
\midrule
Surgery        & 60  & 820  & 5{,}400  & 0.85 \\
Oncology       & 50  & 740  & 4{,}600  & 0.83 \\
Internal Med.  & 40  & 440  & 4{,}300  & 0.82 \\
\midrule
\textbf{Total} & \textbf{150} & \textbf{2{,}000} & \textbf{14{,}300} & \textbf{0.83} \\
\bottomrule
\end{tabular}
}
\caption{Corpus statistics and inter-annotator agreement (IAA).}
\label{tab:medalign_stats}
\end{table}

MedAlign-CN provides a reliable, high-coverage benchmark for evaluating \textit{patient-level temporal reasoning consistency} in Chinese clinical narratives.

\section{Experiments}
\label{sec:exp}

\subsection{Settings and Baselines}
All models are optimized under the unified objective in Eq.~(\ref{eq:train_total}), jointly supervising the generator ($\mathcal{L}_{\text{event}}$), contextual verifier ($\mathcal{L}_{\text{time}}$), and global projection ($\mathcal{L}_{\text{consistency}}$). This ensures every experiment directly evaluates the full \textbf{Generate–Verify–Project (GVP)} framework. We evaluate on the proposed \textbf{MedAlign-CN} corpus (§\ref{sec:dataset}).

\paragraph{Baselines.}
We compare six representative paradigms:
\begin{itemize}[noitemsep,topsep=2pt,leftmargin=1.2em]
  \item \textbf{Span-based IE}~\cite{lin2020joint, wadden2019dygie}: trigger–argument detection without explicit temporal modeling.
  \item \textbf{Nested Event Extraction}~\cite{nagata2020deepeventmine}: handling hierarchical dependencies in biomedical text.
  \item \textbf{Instruction-tuned LLMs}: GPT-4, \texttt{LLaMA-Med} \cite{touvron2023llama}, and \texttt{Baichuan-Med}, prompted with structured JSON templates and Chain-of-Thought (CoT) \cite{wei2022cot}.
  \item \textbf{LLM-RAG-Prompt}: the same FAISS/BM25 hybrid retriever ($k{=}10$) followed by a cross-encoder re-ranker and structured JSON generation. No post-hoc temporal closure or rule repair is applied.
  \item \textbf{Ours (GVP)}: the full differentiable framework.
\end{itemize}
Prompt templates and schemas are in Appendix~\ref{app:prompts}; details in Appendix~\ref{app:impl}.

\paragraph{Backbone and Efficient Fine-tuning.}
We employ \texttt{LLaMA-Med-13B} as the backbone. To enable efficient training on consumer-grade GPUs (e.g., NVIDIA T4 $2{\times}16$GB), we utilize \textbf{4-bit QLoRA} quantization via the \texttt{Unsloth} framework, reducing memory footprint by 60\% while maintaining performance. LoRA adapters (rank $r{=}16$, $\alpha{=}32$) are injected into all linear layers.

\subsection{Evaluation Metrics}
We evaluate:
\begin{itemize}[noitemsep,topsep=2pt,leftmargin=1.2em]
  \item \textbf{Temporal F1}: span-level F1 for events and temporal anchors.
  \item \textbf{TLINK F1}: accuracy of temporal relations.
  \item \textbf{Time-Norm MAE}$\downarrow$: mean absolute error between normalized anchors (in days).
  \item \textbf{Time Extraction Error Rate (TER)$\downarrow$}: proportion of failed or wrongly normalized temporal expressions.
  \item \textbf{Timeline Edit Distance (TED)$\downarrow$}: minimal edit operations to align predicted vs.\ gold timelines.
  \item \textbf{Global Violation Rate (GVR)$\downarrow$}: proportion of graphs violating transitivity or acyclicity; \textbf{Consistency Rate}$\uparrow {=}\,1{-}\text{GVR}$.
  \item \textbf{Expected Calibration Error (ECE)$\downarrow$}: calibration of confidence.
\end{itemize}
Statistical significance is assessed with paired Wilcoxon tests ($p{<}.01$).

\subsection{Main Results}
\begin{table*}[t]
\centering
\setlength{\tabcolsep}{6pt}
\caption{Results on the MedAlign-CN test set (mean ± std over 3 runs). $\downarrow$ lower is better.}
\label{tab:main}
\small
\begin{tabular}{lcccccc}
\toprule
\textbf{Model} & \textbf{Temp-F1} & \textbf{TLINK-F1} & \textbf{MAE$\downarrow$} & \textbf{TER$\downarrow$} & \textbf{TED$\downarrow$} & \textbf{Consist.$\uparrow$} \\
\midrule
Span-based IE      & 71.2 & 58.7 & 6.3 & 12.4\% & 41.3 & 72.5 \\
MRC-style IE       & 74.8 & 60.4 & 5.7 & 11.1\% & 38.6 & 75.3 \\
TIE-Transformer     & 78.0 & 64.9 & 4.9 &  9.2\% & 32.1 & 82.0 \\
GPT-4 (struct.)    & 80.3 & 68.7 & 4.4 &  8.6\% & 29.4 & 84.5 \\
\texttt{LLaMA-Med} & 79.5 & 67.2 & 4.6 &  8.9\% & 30.7 & 83.9 \\
\textbf{LLM-RAG-Prompt} & 81.6 & 70.1 & 4.1 &  7.3\% & 27.6 & 86.7 \\
\midrule
\textbf{Ours (GVP)} & \textbf{84.2±0.3} & \textbf{73.5±0.4} & \textbf{3.2±0.1} & \textbf{5.4±0.2}\% & \textbf{22.4±0.5} & \textbf{92.1±0.2} \\
\bottomrule
\end{tabular}
\end{table*}

The proposed GVP framework achieves the best overall performance, improving Temporal F1 by +5.6 and reducing TED by 32\% over the strongest prior baseline. Crucially, GVP surpasses the \emph{LLM-RAG-Prompt} system, verifying that consistency constraints add value beyond retrieval.

\subsection{Ablation and Analysis}
To quantify the contribution of each module, we perform ablation studies on the MedAlign-CN dev set (Table~\ref{tab:ablation}).

\begin{table}[t]
\centering
\resizebox{\columnwidth}{!}{
\begin{tabular}{lccc}
\toprule
\textbf{Variant} & \textbf{Temp-F1} & \textbf{MAE$\downarrow$} & \textbf{TED$\downarrow$} \\
\midrule
Full (GVP) & \textbf{84.2} & \textbf{3.2} & \textbf{22.4} \\
– w/o $\mathcal{L}_{\text{event}}$ (Schema) & 81.9 & 3.7 & 28.6 \\
– w/o $\mathcal{L}_{\text{time}}$ (Verifier) & 81.0 & 3.9 & 27.8 \\
– w/o $\mathcal{L}_{\text{consistency}}$ (Projection) & 79.6 & 4.6 & 33.9 \\
– w/ Discrete Closure & 80.7 & 4.2 & 31.5 \\
\bottomrule
\end{tabular}
}
\caption{Ablation on the MedAlign-CN dev set.}
\label{tab:ablation}
\end{table}

\paragraph{Impact of Global Projection.}
Removing the differentiable projection layer ($\mathcal{L}_{\text{consistency}}$) causes the largest performance drop in consistency (-12.5\% vs GVP) and a significant increase in Timeline Edit Distance (+11.5). This confirms that without explicit structural regularization, the model struggles to maintain transitivity over long contexts. While discrete closure (row 4) allows for post-hoc repair, it fails to provide training signals, resulting in suboptimal event representations.

\paragraph{Impact of Contextual Verifier.}
The verifier ($\mathcal{L}_{\text{time}}$) is crucial for resolving implicit anchors. Its removal degrades Temporal F1 by 3.2 points, as the model loses the ability to retrieve cross-document evidence for ambiguous time expressions (e.g., "discharge day" vs "post-op day 5").

\paragraph{Schema Constraints.}
Enforcing JSON schema constraints during decoding ($\mathcal{L}_{\text{event}}$) ensures structural validity. Without it, the model occasionally generates malformed trees, leading to a drop in both extraction F1 and overall consistency.

\paragraph{Error Taxonomy.}
Manual inspection of 150 dev-set cases reveals: 63\% of errors stem from ambiguous or implicit anchors (e.g., ``术后恢复期''), 27\% from incomplete context, and 10\% from annotation disagreement. Schema constraints prevent malformed outputs, while projection reduces vague-anchor errors by $\approx$21\%.

\paragraph{Calibration Behavior.}
Structured decoding yields smoother confidence curves and lower ECE across epochs (Figure~\ref{fig:ece_curve}), indicating improved uncertainty calibration under global constraints.

\begin{figure}[t]
\centering
\fbox{
\begin{minipage}[c][38mm][c]{0.9\linewidth}
\small
Placeholder: ECE vs.\ epoch curve showing GVP (solid) vs.\ GPT-4 (dashed) convergence.
\end{minipage}}
\caption{Calibration comparison between GVP and GPT-4.}
\label{fig:ece_curve}
\vspace{-0.5em}
\end{figure}

\subsection{Case Studies}
\label{sec:case}
Figure~\ref{fig:case} shows representative cases. Our model corrects temporal inconsistencies that baselines misinterpret due to implicit Chinese cues.

\begin{figure}[t]
\centering
\fbox{
\begin{minipage}[c][50mm][c]{0.93\linewidth}
\small
\textbf{Example:}\\
\texttt{术后第3天复查CT，未见出血；第5天拔管。}\\
\textit{CT on post-operative day 3 showed no hemorrhage; drain removed on day 5.}\\[0.3em]
\textbf{Baseline (TempTime):}\\
Predicted \textsc{Overlap(CT, Drain Removal)} (incorrect).\\[0.3em]
\textbf{Ours (GVP):}\\
Predicted \textsc{Before(CT, Drain Removal)}, consistent with the globally enforced timeline.
\end{minipage}}
\caption{Case study of temporal reasoning in Chinese clinical text. Implicit markers (“第3天”, “第5天”) often mislead baselines. The projection–verifier synergy restores consistent ordering.}
\label{fig:case}
\vspace{-0.5em}
\end{figure}

\subsection{Extended Analyses}
\label{sec:extended_analysis}

\paragraph{Per-Subset Performance.}
We further analyze performance across clinical specialties (Table~\ref{tab:subset}). \textsc{GVP} consistently improves Temporal F1 by 4–7 points and reduces Timeline Edit Distance (TED) by approximately 50\% across Surgery, Oncology, and Internal Medicine subsets. This robustness indicates that the global consistency constraints generalize well across differing documentation styles.

\begin{table}[h]
\centering
\caption{Per-subset results on MedAlign-CN (test set).}
\label{tab:subset}
\begin{tabular}{lccc}
\toprule
Subset & Temp-F1 & TED$\downarrow$ & Consist.$\uparrow$ \\
\midrule
Surgery        & 84.7 & 21.3 & 91.9 \\
Oncology       & 83.5 & 22.7 & 92.5 \\
Internal Med.  & 82.9 & 23.1 & 91.8 \\
\bottomrule
\end{tabular}
\end{table}

\paragraph{Low-Resource Robustness.}
Training with only 20\% of the labeled data, \textsc{GVP} maintains a Temporal F1 $>$80, whereas baselines drop significantly ($<$72). The schema constraints and global consistency loss act as effective regularizers, preventing overfitting to sparse data.

\paragraph{Error Distribution Analysis.}
We categorize errors into four types: ambiguous anchors (40\%), incomplete context (27\%), fuzzy temporal references (19\%), and annotation disagreement (14\%). The differentiable projection layer is most effective at resolving fuzzy references, reducing their error contribution by $\approx$9 percentage points compared to the ``no projection'' variant.

\section{Discussion}
\label{sec:discussion}

\paragraph{Interpretability and Clinical Alignment.}
The case analyses (\cref{fig:case}) demonstrate that our framework not only improves temporal F1 and coherence metrics but also produces event timelines that align with physicians’ reasoning. By explicitly modeling generation, projection, and verification, the system exposes intermediate decisions that can be audited or corrected—a property rarely achieved by black-box LLMs.

\paragraph{From Local Extraction to Global Reasoning.}
Most temporal IE frameworks focus on local TLINK prediction, which often leads to inconsistencies. Our formulation instead treats temporal extraction as a \textit{global reasoning problem}: events and their temporal anchors are jointly optimized under structural constraints that guarantee transitive closure and acyclicity.

\section{Conclusion}
\label{sec:conclusion}

We have established a \textbf{structure-aware framework} for extracting temporally consistent events from clinical text. By integrating differentiable projection with contextual verification, the model learns to enforce global time coherence internally. Extensive experiments on the new \textsc{MedAlign-CN} benchmark confirm significant improvements in both F1 scores and timeline consistency. We release the dataset and code to facilitate further research into interpretable, logically sound clinical language models.

\section*{Limitations}
The current study focuses on the MedAlign-CN benchmark, which is primarily in Chinese. While the GVP framework is language-agnostic, evaluating its performance on other languages and more diverse clinical settings remains for future work. Additionally, the differentiable projection layer relies on soft constraints; investigating the trade-offs between hard vs. soft logic for absolute temporal consistency is an open question.




\section*{Ethics Statement}
\label{sec:ethics}
\paragraph{Institutional Approval.} The \textbf{MedAlign-CN} corpus was developed under IRB-style oversight. All text data were collected with explicit institutional consent and anonymized.
\paragraph{De-identification.} We adopt a hybrid pipeline combining deterministic regex and transformer-based NER.
\paragraph{Data License.} MedAlign-CN is released under an \textbf{institutional non-commercial research license}.

\bibliography{refs}

\appendix

\section{Annotation Guidelines (Excerpts)}
\label{app:ann_guidelines}
Detailed guidelines for relative time normalization, cross-document linking, and adjudication strategies.

\section{Prompt Templates \& JSON Schema}
\label{app:prompts}

\begin{lstlisting}[language=json, caption={Structured event--time JSON schema.}, label={lst:json-schema}]
{
  "patient_id": "<ID>",
  "events": [
    {
      "trigger": "<span>",
      "arguments": {
        "type": "<Type>",
        "anatomy": "<span?>",
        "value": "<span?>"
      },
      "time": "<YYYY-MM-DD | REL>",
      "evidence": ["<s1>", "<s2>"]
    }
  ],
  "relations": [
    {"head": 0, "tail": 1, "type": "Before|After|Overlap"}
  ]
}
\end{lstlisting}

\begin{lstlisting}[style=prompt, caption={GPT-4 prompt excerpt (English, CN-aware).}, label={lst:llm-baseline-prompt}]
You extract clinical events and times from Chinese notes.
Return ONLY valid JSON following the schema.
Normalize relative times (e.g., Post-op Day 3 -> REL:POSTOP+3D).
\end{lstlisting}

\section{Implementation Details}
\label{app:impl}
For full implementation details, including hyperparameters and training infrastructure, please refer to the source code repository.

\end{CJK*}
\end{document}
